\section{Verification of \texttt{compute} (Tasks 2.1 and 2.2)}
\label{sec:compute}

\subsection{Specification}

\begin{align*}
  \textbf{Precondition:}\quad & \mathit{rev\_sorted}(a, n) \;\wedge\; n \geq 0
    \;\wedge\; \forall\, i < n.\; a[i] \geq 0 \\
  \textbf{Postcondition:}\quad & \hindex(a, n) = h
\end{align*}

\subsection{Loop Invariant}

\begin{invariant}\label{inv:compute}
At the top of each iteration of the \texttt{while} loop in \texttt{compute}:
\[
  I(h) \;\stackrel{\text{def}}{\iff}\;
  0 \leq h \leq n \;\wedge\; \forall\, j.\; 0 \leq j < h \implies a[j] \geq h.
\]
\end{invariant}

\begin{remark}
The invariant states that $h$ is in range and the first $h$ elements of the array
are all at least $h$. Equivalently: $a[j] > j$ for all $j < h$, but the form
$a[j] \geq h$ is stronger and directly useful.
\end{remark}

\subsection{Task 2.1: Safety}

\begin{theorem}[Safety of \texttt{compute}]\label{thm:compute-safety}
Under the precondition, \texttt{compute} terminates and performs no out-of-bounds
array accesses.
\end{theorem}

\begin{proof}
\textbf{Array bounds safety.}
The only array access is $a[h]$ in the loop guard \texttt{h < a[h]}.
Due to short-circuit evaluation, $a[h]$ is evaluated only when $h < n$.
Combined with $h \geq 0$ from the invariant, we have $0 \leq h < n$,
so the access is within bounds.

\textbf{Termination.}
Define the variant function $V(h) = n - h$.
\begin{itemize}
  \item \emph{Bounded below:} By the invariant, $h \leq n$, so $V(h) \geq 0$.
  \item \emph{Strictly decreasing:} Each iteration increments $h$ by~1,
    so $V$ decreases by~1.
\end{itemize}
Since $V$ is a natural number that strictly decreases each iteration,
the loop terminates in at most $n$ iterations.
\end{proof}

\subsection{Task 2.2: Correctness}

\begin{theorem}[Correctness of \texttt{compute}]\label{thm:compute-correct}
Under the precondition, \texttt{compute} returns $\hindex(a, n)$.
\end{theorem}

\begin{proof}
We prove the loop invariant by induction, then derive the postcondition.

\medskip\noindent
\textbf{Initialization.}
Before the loop, $h = 0$.
\begin{itemize}
  \item $0 \leq 0 \leq n$: holds since $n \geq 0$.
  \item $\forall\, j.\; 0 \leq j < 0 \implies a[j] \geq 0$: vacuously true.
\end{itemize}
Thus $I(0)$ holds.

\medskip\noindent
\textbf{Maintenance.}
Assume $I(h)$ holds and the guard is true: $h < n \wedge h < a[h]$.
After executing $h \mathrel{{+}{+}}$, the new value is $h' = h + 1$. We show $I(h')$:
\begin{itemize}
  \item \emph{Range:} $h' = h + 1 \leq n$ since $h < n$ (guard) and $h, n \in \Z$.
    Also $h' = h + 1 \geq 1 > 0$.
  \item \emph{Universal property:} We must show $\forall\, j.\; 0 \leq j < h' \implies a[j] \geq h'$,
    i.e., $\forall\, j.\; 0 \leq j < h + 1 \implies a[j] \geq h + 1$.

    \textbf{Case $j < h$:}
    By $I(h)$, $a[j] \geq h$.
    Since the array is reverse-sorted and $j < h$, we have $a[j] \geq a[h]$.
    From the guard, $a[h] > h$, so $a[h] \geq h + 1$ (integer arithmetic).
    Therefore $a[j] \geq a[h] \geq h + 1 = h'$.

    \textbf{Case $j = h$:}
    From the guard, $a[h] > h$, so $a[h] \geq h + 1 = h'$.
\end{itemize}
Thus $I(h + 1) = I(h')$ holds.

\medskip\noindent
\textbf{Postcondition.}
When the loop exits, $I(h)$ holds and the guard is false:
$h \geq n$ or $a[h] \leq h$.
Since $I(h)$ gives $h \leq n$, the disjunction becomes: $h = n$ or ($h < n$ and $a[h] \leq h$).

By Corollary~\ref{cor:transition-equiv}, $\hindex(a,n) = h$ if and only if
$(\forall\, j < h.\; a[j] > j)$ and ($h = n$ or $a[h] \leq h$).

The second condition holds by the negated guard. For the first condition:
from $I(h)$, for all $j < h$: $a[j] \geq h > j$ (since $j < h$), giving $a[j] > j$.

Therefore $\hindex(a,n) = h$, and \texttt{compute} returns the h-index.
\end{proof}

\begin{remark}[Alternative invariant formulation]
An equivalent invariant is:
$I'(h) \iff 0 \leq h \leq n \wedge \forall\, j < h.\; a[j] > j$.
This follows from $I(h)$ since $a[j] \geq h > j$ for $j < h$.
Conversely, $a[j] > j$ with reverse-sorting gives $a[j] \geq a[h{-}1] > h{-}1$,
so $a[j] \geq h$. The formulation $a[j] \geq h$ is stronger and simplifies
the postcondition argument.
\end{remark}
